\chapter{General Introduction}
\label{ch:introduction}

\vcyclephase{Project Initiation --- Problem Definition}

\section{Context of the Internship}
\label{sec:context}

This report presents the work carried out during the final-year engineering project (Projet de Fin d'\'Etudes --- PFE) within the framework of the mechatronics engineering curriculum. The project focuses on the design and implementation of a \textbf{Smart Storage Management System (SSMS)} --- a Pick/Put-to-Light warehouse bin guidance system aimed at improving the efficiency and accuracy of manual warehouse operations.

\section{Problem Statement}
\label{sec:problem}

In traditional manual warehouses, operators locate storage bins by reading labels, paper pick lists, or handheld terminal displays. This process is prone to several inefficiencies:
\begin{itemize}
  \item \textbf{Bin-finding errors:} Operators may misread labels or navigate to the wrong location, leading to picking errors that propagate through the supply chain.
  \item \textbf{Slow picking speed:} Without visual guidance, each pick requires visual search time, reducing overall throughput.
  \item \textbf{No real-time feedback:} Errors are typically discovered only at downstream verification stations or after shipment, causing costly returns and customer dissatisfaction.
  \item \textbf{Lack of traceability:} Manual processes offer limited data collection on operator performance and picking accuracy.
\end{itemize}

\section{Proposed Solution: SSMS}
\label{sec:solution}

The Smart Storage Management System addresses these challenges by providing real-time visual guidance to warehouse operators using addressable RGB LED strips mounted at each storage bin. When an item needs to be picked or stored, the corresponding bin illuminates with a specific colour and pattern, guiding the operator directly to the correct location. The system consists of:

\begin{itemize}
  \item An \textbf{ESP32-WROOM-32D}\cite{esp32_trm} microcontroller running MicroPython\cite{micropython_docs} firmware that controls the LED strips based on MQTT commands.
  \item \textbf{WS2813 addressable RGB LEDs}\cite{ws2813_ds} (6 LEDs per bin, 36 bins, 216 total) providing bright, unambiguous visual cues.
  \item \textbf{74HCT-series level-shifting ICs} (SN74HCT244, SN74HCT245, SN74HCT573) to interface the 3.3\,V ESP32 outputs with the 5\,V WS2813 data line.
  \item A \textbf{Python/Tkinter desktop application} running on a Windows PC that operators interact with via barcode scanning.
  \item \textbf{MQTT communication}\cite{oasis_mqtt} over a local LAN using a Mosquitto broker for reliable messaging between the desktop application and the ESP32 controller.
\end{itemize}

\section{Report Structure}
\label{sec:structure}

This report follows the \textbf{V-Cycle (Cycle en V)} methodology, a structured systems engineering approach well-suited to embedded mechatronic projects. The document is organised as follows:

\begin{description}
  \item[Chapter 1 --- General Context] introduces the operational environment, project objectives, stakeholders, and functional analysis diagrams (B\^ete \`a Cornes and Pieuvre/Octopus).
  \item[Chapter 2 --- Requirements Specification] documents the functional and non-functional requirements, hardware and software constraints, and communication protocol specifications.
  \item[Chapter 3 --- Solution Design] presents the system architecture, detailed hardware design (schematic, BOM, level-shifting), firmware architecture (state machine, MQTT parsing), desktop application design (MVC pattern, workflow), and simulation strategies (LTspice, Proteus).
  \item[Chapter 4 --- Implementation] describes the physical realisation, firmware and software code with key excerpts, simulation results, integration testing, and requirements traceability.
  \item[Chapter 5 --- General Conclusion] summarises achievements, discusses difficulties encountered, identifies limitations, and proposes future work.
  \item[Chapter 6 --- Annexes] provides supplementary materials: full BOM, complete source code listings, MQTT reference table, schematic, netlist, acronyms, and references.
\end{description}

\section{V-Cycle Methodology Justification}
\label{sec:vcycle}

The V-Cycle methodology is particularly well-suited to this project for several reasons:

\begin{itemize}
  \item \textbf{Embedded system nature:} SSMS involves both hardware (microcontroller, LEDs, level shifters) and software (firmware, desktop application), requiring a structured approach that handles both domains.
  \item \textbf{Validation-driven:} The V-Cycle emphasises verification and validation at each stage, ensuring that requirements trace through to final testing --- critical for a system where reliability directly impacts warehouse operations.
  \item \textbf{Iterative refinement:} The left side of the V (requirements $\rightarrow$ design) allows thorough specification before implementation, while the right side (testing $\rightarrow$ validation) provides systematic verification.
  \item \textbf{Mechatronic alignment:} The V-Cycle is widely adopted in mechatronics and automotive embedded systems, making it the natural choice for this interdisciplinary project.
\end{itemize}

\begin{figure}[H]
  \centering
  \begin{tikzpicture}[
      node distance=1.8cm,
      phase/.style={rectangle, draw, minimum width=3.5cm, minimum height=1cm, align=center, font=\small\bfseries, fill=ssmsBlue!10},
      arrow/.style={-{Latex[length=3mm]}, thick},
      label/.style={font=\small\itshape}
    ]
    \node[phase] (req) at (0,0) {Requirements\\Analysis};
    \node[phase] (val) at (8,0) {Validation\\\& Acceptance};
    \node[phase] (func) at (1.5,-1.8) {Functional\\Design};
    \node[phase] (integ) at (6.5,-1.8) {Integration\\Testing};
    \node[phase] (det) at (3,-3.6) {Detailed\\Design};
    \node[phase] (utest) at (5,-3.6) {Unit\\Testing};
    \node[phase] (impl) at (4,-5.4) {Implementation};

    \draw[arrow] (req) -- node[above,label] {specification} (val);
    \draw[arrow] (req) -- (func);
    \draw[arrow] (val) -- (integ);
    \draw[arrow] (func) -- (det);
    \draw[arrow] (integ) -- (utest);
    \draw[arrow] (det) -- (impl);
    \draw[arrow] (impl) -- (utest);

    \node[above=0.3cm, align=center, font=\bfseries] at (4,0.5) {V-Cycle Methodology Applied to SSMS};
  \end{tikzpicture}
  \caption{V-Cycle methodology applied to the SSMS project.}
  \label{fig:vcycle}
\end{figure}
